\section{Service Deployment}

\begin{frame}{Section 6.4: Service Deployment}
\begin{itemize}
    \item After a system is built, it must be deployed, monitored, and updated.
    \item In microservices systems, deployment is more complex than in monolithic systems.
    \item Different services may use different languages, libraries, databases, and runtime environments.
    \item This makes deployment a major architectural and operational concern.
\end{itemize}
\end{frame}

\begin{frame}{Why Deployment Is Harder with Microservices}
\begin{itemize}
    \item A microservices system may contain tens or hundreds of services.
    \item Each service may have:
    \begin{itemize}
        \item its own dependencies,
        \item its own runtime requirements,
        \item its own update cycle.
    \end{itemize}
    \item There is no single standard deployment configuration for all services.
    \item A separate operations team can easily become a deployment bottleneck.
\end{itemize}
\end{frame}

\begin{frame}{From Separate Ops to DevOps}
\begin{itemize}
    \item Traditionally, operations and development were separate activities.
    \item In microservices systems, it is common for development teams to take responsibility for:
    \begin{itemize}
        \item deployment,
        \item service management,
        \item ongoing operational support.
    \end{itemize}
    \item This combined approach is commonly called \textbf{DevOps}.
\end{itemize}
\end{frame}

\begin{frame}{Continuous Deployment}
\begin{itemize}
    \item A general principle of microservice-based development is that teams own deployment of their service.
    \item A common practice is \textbf{continuous deployment}.
    \item When a change is made and validated, the modified service is re-deployed automatically.
    \item This contrasts with older release models where changes were bundled into occasional large releases.
\end{itemize}
\end{frame}

\begin{frame}{A Continuous Deployment Pipeline}
\begin{itemize}
    \item A typical pipeline includes:
    \begin{itemize}
        \item commit code changes,
        \item run unit tests,
        \item build the test system,
        \item run integration tests,
        \item package the service,
        \item deploy the new version,
        \item run acceptance or whole-system tests,
        \item replace the current service if all tests pass.
    \end{itemize}
    \item If a step fails, the change is rejected.
\end{itemize}
\end{frame}

\begin{frame}{Why Automation Matters}
\begin{itemize}
    \item Continuous deployment depends on automation.
    \item Automated tests must check that:
    \begin{itemize}
        \item the modified service works correctly,
        \item it still interacts correctly with other services,
        \item the whole system remains stable after deployment.
    \end{itemize}
    \item Without automation, rapid deployment quickly becomes dangerous chaos.
\end{itemize}
\end{frame}

\begin{frame}{Containers as Deployment Units}
\begin{itemize}
    \item Containers are often the preferred way to package microservices for deployment.
    \item A container includes the executable service and the software environment it needs.
    \item This makes deployment more portable across servers.
    \item Containers reduce the need for service teams to worry about low-level server configuration differences.
\end{itemize}
\end{frame}

\begin{frame}{Why Containers Help}
\begin{itemize}
    \item Containers allow dependencies to be defined in advance.
    \item A service can often be deployed by:
    \begin{itemize}
        \item building the container image,
        \item loading it onto a server or cloud platform,
        \item starting the containerized service.
    \end{itemize}
    \item This makes deployment more repeatable and easier to automate.
\end{itemize}
\end{frame}

\begin{frame}{Managing Many Containers}
\begin{itemize}
    \item Large microservice systems may involve tens, hundreds, or thousands of containers.
    \item Managing them manually is difficult.
    \item Container-management systems help automate:
    \begin{itemize}
        \item deployment,
        \item scheduling,
        \item service discovery,
        \item load balancing,
        \item resource management.
    \end{itemize}
\end{itemize}
\end{frame}

\begin{frame}{Kubernetes and Container Orchestration}
\begin{itemize}
    \item The book highlights \textbf{Kubernetes} as an example of a container-management platform.
    \item Kubernetes helps manage container deployment across clusters of servers.
    \item It supports:
    \begin{itemize}
        \item automated deployment,
        \item scheduling,
        \item service discovery,
        \item load balancing.
    \end{itemize}
    \item This kind of orchestration becomes essential at scale.
\end{itemize}
\end{frame}

\begin{frame}{Deployment Risk}
\begin{itemize}
    \item A major deployment risk is that a new service version may interact badly with existing services.
    \item Even when tests pass, unexpected behavior may still appear in production.
    \item Testing reduces risk, but it cannot eliminate it completely.
    \item Therefore, deployed services must be monitored carefully after release.
\end{itemize}
\end{frame}

\begin{frame}{Monitoring New Versions}
\begin{itemize}
    \item After deployment, the system should watch for:
    \begin{itemize}
        \item failures,
        \item degraded performance,
        \item unexpected interactions,
        \item incorrect responses.
    \end{itemize}
    \item Monitoring helps detect whether a new version is safe to keep in production.
    \item If serious problems appear, the system should be able to roll back.
\end{itemize}
\end{frame}

\begin{frame}{Rollback and Version Switching}
\begin{itemize}
    \item One practical strategy is to keep the old version available while introducing the new version.
    \item If the monitoring system detects problems, traffic can be switched back to the older version.
    \item This makes rollback faster and safer than rebuilding the old version from scratch.
    \item API gateways can help route requests to the currently active service version.
\end{itemize}
\end{frame}

\begin{frame}{Versioned Services}
\begin{itemize}
    \item In some systems, service versions are identified explicitly.
    \item For example, an API change may require clients or gateways to distinguish between:
    \begin{itemize}
        \item version 001,
        \item version 002,
        \item and later versions.
    \end{itemize}
    \item Over time, clients and gateways can be updated to use the newer service version.
\end{itemize}
\end{frame}

\begin{frame}{Big Picture: What Deployment Requires}
\begin{itemize}
    \item Effective microservice deployment depends on:
    \begin{itemize}
        \item team ownership,
        \item automation,
        \item testing,
        \item containers,
        \item monitoring,
        \item safe rollback mechanisms.
    \end{itemize}
    \item Deployment is not an afterthought.
    \item It is part of the architecture.
\end{itemize}
\end{frame}

\begin{frame}{Section 6.4 Recap}
\begin{itemize}
    \item Microservice deployment is more complex than monolithic deployment because services are independently developed and updated.
    \item DevOps and continuous deployment help teams manage this complexity.
    \item Containers provide a portable deployment unit, and Kubernetes helps manage containers at scale.
    \item Monitoring and rollback are essential because even tested deployments can fail in production.
\end{itemize}
\end{frame}

\begin{frame}{Looking Ahead}
\centering
\Large Next: Summary and Further Reading

\vspace{1em}

\normalsize
We now move from chapter content into final takeaways and recommended reading.
\end{frame}
